\documentclass[11pt,a4paper]{letter}
\usepackage[utf8]{inputenc}
\usepackage[T1]{fontenc}
\usepackage[french]{babel}
\usepackage[left=2.5cm,right=2.5cm,top=2cm,bottom=2cm]{geometry}
\usepackage{hyperref}
\usepackage{xcolor}

\definecolor{primary}{RGB}{35, 55, 123}
\hypersetup{colorlinks=true, urlcolor=primary}

\signature{Loïc Aron MBASSI EWOLO}
\address{Loïc Aron MBASSI EWOLO\\Cergy, France\\(+33) 7 59 52 33 98\\loicaron.mbassiewolo@ensea.fr}

\begin{document}

\begin{letter}{Mme Noemi LANCIOTTI\\Laboratoire GeePs\\CentraleSupélec\\3 rue Joliot Curie\\91190 Gif-sur-Yvette}

\opening{Madame,}

Actuellement élève-ingénieur en deuxième année à l'ENSEA dans le cadre d'un double diplôme avec l'École Polytechnique de Yaoundé, je me permets de vous adresser ma candidature pour le stage « Caractérisation des performances du convertisseur triphasé open-source de la suite OwnTech ».

Ce sujet correspond parfaitement à mes aspirations : allier électronique de puissance, expérimentation en laboratoire et contribution à un projet open-source. La démarche scientifique portée par OwnTech, favorisant la reproductibilité et le partage des résultats, correspond exactement à ma vision de l'ingénierie.

Ma formation à l'ENSEA m'apporte des bases solides en électronique analogique et numérique, en conception FPGA (VHDL) et en automatique. Parallèlement, mes projets personnels m'ont permis de développer des compétences directement applicables à ce stage :

\begin{itemize}
    \item \textbf{Prototypage hardware :} Au sein du laboratoire IOTA de Polytechnique, j'ai participé au projet Harmony Gloves, où j'ai réalisé la soudure de composants sur PCB, intégré des capteurs (IMU, flex sensors) sur microcontrôleur STM32, et mené des campagnes d'acquisition de données expérimentales.
    
    \item \textbf{Systèmes embarqués :} Dans le cadre du challenge CoHoMa (robotique de défense), je travaille actuellement sur l'optimisation d'algorithmes sur Nvidia Jetson, avec intégration de capteurs (LiDAR 3D, caméras) et tests expérimentaux en conditions réelles.
    
    \item \textbf{Programmation embarquée :} Je maîtrise le langage C (gestion mémoire, IPC, temps réel) et j'ai une expérience en VHDL à travers mes projets FPGA actuels à l'ENSEA.
    
    \item \textbf{Instrumentation et traitement de données :} Mes travaux de recherche au MILA et mes projets académiques m'ont habitué à la rigueur des protocoles expérimentaux, au traitement de données sous Python (NumPy, Pandas, Matplotlib) et à la documentation technique.
    
    \item \textbf{Open-source :} Je suis l'auteur d'un package open-source publié sur Packagist (30+ étoiles GitHub), ce qui m'a familiarisé avec les bonnes pratiques de documentation et de partage de code.
\end{itemize}

Je suis particulièrement motivé par la perspective de mettre en place un banc de test, de concevoir les procédures logicielles embarquées et de réaliser une campagne d'essais automatisée. La valorisation des travaux sur docs.OwnTech.org correspond à mon goût pour la diffusion scientifique et technique.

Disponible dès avril 2026 pour une durée de 5 mois, je serais ravi d'échanger avec vous sur ce sujet passionnant.

\closing{Dans l'attente de votre retour, je vous prie d'agréer, Madame, l'expression de mes salutations distinguées.}

\end{letter}
\end{document}
